\paragraph{Numerical evidence for Zipf's and Heaps' laws.}
To test whether the prime product lattice supports the same statistical
signatures as Loreto et al.'s P\'olya urn, we performed
500~independent random walks of 10{,}000~steps on the Boolean lattice
of square-free products of the first~20~primes
($2^{20} = 1{,}048{,}576$~nodes).  At each step a walker selects
prime~$p_i$ with probability proportional to~$n_i + \rho$, where $n_i$
counts prior selections of~$p_i$ and $\rho > 0$ controls the
exploration--exploitation trade-off.  We tested two dynamics---a
\emph{toggle} model (the selected prime is added if absent, removed if
present) and an \emph{additive} model (primes accumulate
irreversibly)---sweeping $\rho \in \{0.1, 0.5, 1.0, 5.0\}$.

In the toggle model with strong reinforcement ($\rho = 0.1$), the
rank-frequency distribution follows a power law with exponent
$\alpha = 0.76$ ($R^2 = 0.978$, fit to the top quartile of ranks),
and Heaps' law holds with $\beta = 0.81$ ($D(t) \sim t^{0.81}$).
As~$\rho$ increases (weaker reinforcement), the walk explores the
lattice nearly uniformly and the Zipf exponent collapses toward zero.

The additive model yields the most striking result.  At $\rho = 5.0$,
the rank-frequency distribution is an almost exact Zipf law:
$\alpha = 1.00$ ($R^2 = 0.97$) over the full rank range.  Intermediate
values produce steeper power laws ($\alpha = 1.47$ at $\rho = 1.0$,
$\alpha = 2.01$ at $\rho = 0.5$), tracing a continuous transition from
steep to Zipfian scaling as the novelty bias grows.  The additive model's
irreversibility concentrates visits on a small number of lattice
trajectories ($\sim 7{,}600$~unique nodes), creating the
frequency heterogeneity that underlies Zipf's law.

These results demonstrate that P\'olya-like reinforcement on the prime
product lattice can reproduce the Zipf and Heaps signatures of
\citet{loreto2017}'s urn model.  The additive dynamics, which most
closely mirror irreversible discovery, produce canonical Zipf scaling
($\alpha \approx 1$) when the novelty parameter is appropriately tuned.
The exponent's continuous dependence on~$\rho$ offers a route to
calibrate the lattice model against empirical data from different
discovery domains.  This confirms that the lattice provides a valid
structural complement to the stochastic urn framework: it not only
captures the \emph{structure} of the adjacent possible (as shown in
Sections~3--4) but, when coupled with reinforced exploration, also
reproduces its \emph{statistical} signatures.
